\chapter{Two Axes: Recognition and Implementation}
\label{ch:two-axes}

\begin{plainlanguage}
\textbf{In plain terms:} Knowing the right answer and actually doing it are not the same skill. This chapter formally separates \emph{recognition} (telling right from wrong) from \emph{implementation} (acting on it once recognized), and shows that most of what gets dismissed as foolishness in the Kalabari examples is really an implementation failure, not a recognition one.
\end{plainlanguage}

\section{A single margin obscures two different difficulties}

The Kalabari material resists a single-margin story, and resisting it is what makes the material useful rather than merely illustrative. Consider two examples from the childhood (\emph{awome}) stage of Wariboko's account. The bucket-and-basket test is difficult, if at all, only for a child who has not yet grasped basic physical affordances --- once the affordance is understood, the choice is trivial, because nothing pulls the child toward the basket. Compare this to the drink-sharing ritual reported in the same source: the youngest male in a gathering is expected to serve the elders first and take only what remains, even when little is left and he plainly wants more. Here the difficulty is not perceptual at all. The child fully understands the rule. The challenge is resisting the pull of his own thirst.

These are different problems, and collapsing them into a single scalar margin obscures the difference. It is more accurate to track two quantities separately.

\begin{definition}[Two-Axis Decomposition]
Let $d_p(a,c)$ denote the perceptual or conceptual distinguishability of $a$ from $c$ --- how obvious the boundary is. Let $m(a,c)$ denote the motivational pull of the inadmissible option $c$, independent of how obvious the boundary is.
\end{definition}

The bucket-basket case is high-$d_p$, near-zero-$m$: the boundary is glaring, and nothing tempts the child toward the wrong side of it. The drink-sharing case is near-zero-$d_p$, high-$m$: the boundary is perfectly clear, and the difficulty lies entirely in enacting it against one's own desire. The bibife marriage feast --- offered to a grown bride who is, by the ethnographer's own account, genuinely tempted by the aroma of food she knows is rightfully hers --- is the adult-life extremity of the same axis: $d_p$ effectively zero, $m$ high and sustained over an entire ceremony.

\begin{definition}[Recognition and Implementation Competence]
\[
I_R = P(f(x) = y)
\]
--- the probability of correctly classifying a candidate option as admissible or inadmissible ---
\[
I_I = P(\pi(x) = a \mid f(x) = a)
\]
--- the probability of correctly enacting the recognized-admissible choice, conditional on having recognized it correctly in the first place.
\end{definition}

The conditional form of $I_I$ is the important move here: by defining implementation success only over cases where recognition has already succeeded, the two quantities are cleanly separated by construction. A failure can now be diagnosed as belonging to one or the other, rather than collapsing into a single undifferentiated ``made the wrong choice.'' Much of what gets called stupidity or failure across the Kalabari life-cycle examples is not a recognition failure at all. The bride knows exactly what she wants and exactly what is expected of her. The chief at his installation, asked to choose between a yam (self-provision) and a cannon ball (public defense), is not confused about which item the community demands he select. The soldiers who refuse to catapult their general are not uncertain about what their commander has, in his battle-fury, miscalculated. In each case recognition is not the bottleneck. Execution is.

\begin{proposition}[Dissociation]
High $I_R$ does not imply high $I_I$.
\end{proposition}

This is not given a formal proof here; it is supported constructively, by exhibiting cases --- the adult-stage examples above --- in which $I_R$ is evidently near $1$ while $I_I$ is genuinely under strain. A fuller, independently-grounded version of this proposition is developed in Chapter~\ref{ch:transfer}, once the relevant psychological literature has been brought in.

\subsection*{Four failure modes}

The conditional definition of $I_I$ in terms of $I_R$ makes it possible to lay out, exhaustively, every combination of high and low standing on the two axes, which gives the later transfer matrix of Chapter~\ref{ch:transfer} an intuitive grounding before any matrix notation appears.

\begin{center}
\begin{tabular}{c c p{7.5cm}}
\hline
\textbf{Recognition} & \textbf{Implementation} & \textbf{Interpretation} \\
\hline
Low & Low & ignorance --- neither the boundary nor its enactment is in place \\[2pt]
High & Low & weakness --- the boundary is seen clearly but not held to \\[2pt]
Low & High & luck or habit --- correct behavior without underlying understanding \\[2pt]
High & High & competence --- both recognized and reliably enacted \\
\hline
\end{tabular}
\end{center}

The middle two rows are the cells this book's recognition/implementation split was built to make visible, since a single undifferentiated notion of ``success'' or ``failure'' collapses them into indistinguishable outcomes. \emph{Weakness} is the cell most of this chapter's Kalabari examples occupy --- the bride, the chief, the soldiers all sit in this row, not the ignorance row a flatter theory of failure would place them in. \emph{Luck or habit} is the more unsettling cell, easy to overlook: a person or system can behave correctly for an extended run without the recognition component ever having been present, which means correct behavior alone, observed from outside, cannot by itself distinguish this row from competence --- a methodological caution worth carrying forward into Chapter~\ref{ch:transfer}'s discussion of why transfer, rather than performance alone, is the test that actually discriminates between them.

\subsection*{A fifth case the table cannot represent}

The table above assumes, without stating it, that $I_R$ and $I_I$ belong to the same agent being evaluated --- that whatever recognized the boundary and whatever enacted it are the same locus of competence, even when they come apart from each other as the weakness and luck-or-habit rows describe. There is a fifth situation this assumption hides, worth naming explicitly rather than leaving for a reader to discover as an apparent exception to the table.

Shannon Vallor, in the panel discussion already cited in Chapter~\ref{ch:room-for-river}, describes a four-stage progression she terms \emph{cognitive offloading}, \emph{cognitive debt}, \emph{cognitive deskilling}, and finally \emph{cognitive surrender}: a sequence in which a person first lets an external system perform cognitive work on their behalf, then accumulates a growing gap between the work they could do unassisted and the work they have actually been practicing, then loses the underlying capacity such that performance degrades once the assisting system is removed, and finally arrives at a state of not being aware that the delegation has occurred at all --- continuing to experience themselves as the one deciding, while the actual recognition and implementation have migrated elsewhere.

This describes a competence credited to an agent in the high-$I_R$, high-$I_I$ cell of the table above, observed from outside, that does not belong there once the locus of recognition and implementation is correctly identified --- because the agent credited with the competence is not, by the final stage of Vallor's cascade, the agent who actually possesses it. This is not the same as the luck-or-habit row, where correct behavior occurs without recognition ever having been present in anyone; here recognition and implementation are both genuinely present, just not in the party the outcome gets attributed to. Call this \emph{delegated competence}, distinct from all four cells above precisely because it requires relaxing the table's hidden single-agent assumption rather than moving within it. The unsettling part of Vallor's evidence, on her account, is not that this happens to novices or to populations conventionally treated as vulnerable, but that two independent surveys found it occurring in business leadership specifically --- roughly half of surveyed chief executives, across both a UK and a US sample, reporting that they deferred to an AI system's conclusion over their own judgment when the two conflicted, a population an observer would, by performance alone, have every reason to place in the competence cell.
